% Generated draft: requires mathematical and source review.
\documentclass[11pt]{article}
\usepackage[margin=25mm]{geometry}
\usepackage{fontspec}
\setmainfont{Latin Modern Roman}
\usepackage{amsmath,amssymb,mathtools,unicode-math}
\setmathfont{Latin Modern Math}
\usepackage{enumitem,needspace,xurl,hyperref}
\setlength{\parindent}{0pt}
\setlength{\parskip}{5pt}
\setlength{\emergencystretch}{3em}
\newcommand{\N}{\mathbb{N}}
\newcommand{\Z}{\mathbb{Z}}
\newcommand{\Q}{\mathbb{Q}}
\newcommand{\R}{\mathbb{R}}
\newcommand{\C}{\mathbb{C}}
\newcommand{\F}{\mathbb{F}}
\newcommand{\A}{\mathbb{A}}
\newcommand{\PP}{\mathbb{P}}
\newcommand{\E}{\mathbb{E}}
\newcommand{\eps}{\varepsilon}
\newcommand{\dd}{\,\mathrm{d}}
\begin{document}
\section*{UnsolvedMath ID: 30004939}
\textbf{Source problem code:} catalogue entry
\subsection{Definitions and mathematical statement}
For an integer \(k\geq 0\), let \(S^k=\{x\in\R^{k+1}:\|x\|=1\}\). Write \(X\cong Y\) for homeomorphism and \(X\simeq Y\) for homotopy equivalence. Products of spheres carry the product topology. For a cluster algebra \(\mathcal A\), its cluster complex \(\Delta(\mathcal A)\) is the abstract simplicial complex whose vertices are cluster variables and whose faces are finite subsets contained in a cluster; topological assertions refer to its geometric realization \(|\Delta(\mathcal A)|\). The rank is the number of mutable variables in a seed.

Let \(\operatorname{Ass}_{\mathrm{aff}}(\mathcal A)\) and \(\operatorname{Ass}_{\mathrm{dbl}}(\mathcal A)\) denote, respectively, the affine generalized associahedron and the doubly extended associahedron intended in the cited report. These are placeholders for the report's constructions: the supplied record does not define their cells, identifications, or topology. It also does not define the precise class of doubly extended cluster algebras or specify coefficient conventions.

The supplied passage attributes the following assertions to Conjecture 7:
\begin{enumerate}
\item For every affine cluster algebra \(\mathcal A\) of rank \(n+1\),
\[
\operatorname{Ass}_{\mathrm{aff}}(\mathcal A)\cong S^n.
\]
\item For every doubly extended cluster algebra \(\mathcal A\) of rank \(n+2\),
\[
|\Delta(\mathcal A)|\simeq S^{n-1}.
\]
\item For every doubly extended cluster algebra \(\mathcal A\), the associated doubly extended associahedron is asserted to have a product-of-spheres description. The apparent intended formulas are
\[
\operatorname{Ass}_{\mathrm{dbl}}(\mathcal A)\cong
\begin{cases}
S^{n-1}\times S^2,&\text{outside the exceptional type},\\
S^7\times S^1\times S^1,&\text{in the exceptional type}.
\end{cases}
\]
The exceptional type appears to be \(E_8^{(1,1)}\), but its superscript is displaced in the supplied extraction. The exponents in the exceptional sphere product are likewise displaced. Thus this display is an explicitly tentative reconstruction, not a verified transcription. The use of \(\operatorname{rank}(\mathcal A)=n+2\) in this third assertion appears to continue the convention in the second assertion, but is not restated in the extracted third assertion.
\end{enumerate}
The assertions are universally quantified over the source's intended classes. Their precise admissible ranks are not supplied. The sphere notation defined above covers the first assertion when \(n\geq0\) and the other assertions when \(n\geq1\); the record does not resolve any intended lower-rank cases or convention for \(S^{-1}\). These missing construction, scope, and transcription details must be resolved before this becomes a fully specified mathematical statement.
\subsection{Short English statement}
Conjecture 7 predicts that affine generalized associahedra are spheres, that doubly extended cluster complexes have the homotopy type of spheres, and that doubly extended associahedra are products of spheres with one exceptional type. The source constructions, rank bounds, and damaged exceptional formula require clarification. This is a statement-only draft requiring human mathematical review.
\subsection{Sources}
\begin{itemize}
\item Supplied local UnsolvedMath registry record, permanent ID \(30004939\), title \emph{Topology of Affine and Doubly Extended Associahedra}, source problem code \texttt{catalogue entry}, in the Oberwolfach Reports Open Problems collection. The record identifies the passage as Conjecture 7; no individual proposer is supplied.
\item \emph{Mini-Workshop: Scattering Amplitudes, Cluster Algebras, and Positive Geometries} (2022), Oberwolfach Reports, pp.\ 3101--3133, DOI \texttt{10.4171/owr/2021/57}: \url{https://doi.org/10.4171/owr/2021/57}. This is the authoritative report citation supplied by the record; its contents have not been independently consulted for this draft. The registry separately gives proposed year 2021 and an \texttt{owr-year:2021} tag.
\item Supplied literature source, labeled \emph{cluster modular groups of affine and doubly extended cluster algebras}: \url{https://arxiv.org/abs/2107.10334}.
\item Supplied literature source, labeled \emph{current dissertation statement}: \url{https://api.drum.lib.umd.edu/server/api/core/bitstreams/b8a1780c-1d92-4a30-b06f-56ce0017d770/content}. No author, dissertation title, or publication date is supplied.
\item Literature assessment dated exactly \texttt{2026-08-22}: the supplied curation labels the problem open, reports that the three assertions remain conjectures in subsequent research, and says that no proof covering all listed types was located. This is supplied dated curation, not an independently verified literature assessment. The record explicitly requests independent verification of time-sensitive status and cited claims.
\item The registry labels the statement \texttt{corrected\_verified} and reports that the official Oberwolfach report was checked directly and that the clean statement restores definitions and removes extraction artifacts or source-dependent wording. These are supplied verification claims, not verification performed for this draft; the supplied clean statement does not itself provide the missing constructions or exceptional formula.
\end{itemize}
\end{document}
